\documentclass[12pt,a4paper]{article}

% Page setup and margins
\usepackage[top=1in, bottom=1in, left=1in, right=1in]{geometry}

% Font packages
\usepackage{times}
\usepackage[T1]{fontenc}

% Essential packages
\usepackage{graphicx}
\usepackage{array}
\usepackage{tabularx}
\usepackage{longtable}
\usepackage{multirow}
\usepackage{multicol}
\usepackage{booktabs}
\usepackage{amsmath}
\usepackage{enumitem}
\usepackage{float}
\usepackage{setspace}
\usepackage{titlesec}
\usepackage{fancyhdr}
\usepackage[hidelinks]{hyperref}
\usepackage{url}
\usepackage{amssymb}

% Set single spacing globally
\singlespacing

% Define custom section formatting
\titleformat{\section}
  {\fontsize{16}{19.2}\selectfont\bfseries}
  {\thesection}
  {1em}
  {}

\titleformat{\subsection}
  {\fontsize{13}{15.6}\selectfont\bfseries}
  {\thesubsection}
  {1em}
  {}

\titleformat{\subsubsection}
  {\fontsize{12}{14.4}\selectfont\bfseries\itshape}
  {\thesubsubsection}
  {1em}
  {}

% Remove page numbers from header, add to footer
\pagestyle{fancy}
\fancyhf{}
\fancyfoot[C]{\thepage}
\renewcommand{\headrulewidth}{0pt}
\renewcommand{\footrulewidth}{0pt}

% Custom commands
\newcommand{\checkbox}{\makebox[0pt][l]{$\square$}\raisebox{.15ex}{\hspace{0.1em}$\square$}}
\newcommand{\checkedbox}{\makebox[0pt][l]{$\square$}\raisebox{.15ex}{\hspace{0.1em}$\checkmark$}}

\begin{document}

% Title page
\begin{center}
{\Large \textbf{CUJ-to-Database Agent}}

\vspace{0.5cm}

{\Large \textbf{An LLM-Driven Multi-Agent System for Automated Database Engineering}}

\vspace{0.5cm}

{\Large \textbf{Final Year Project Proposal}}

\vspace{0.5cm}

{\large Session Spring 2026}

\vspace{0.5cm}

{\large 4th Year Student}

\vspace{1cm}

{\large BSc. (Hons.) in Computer Science / Artificial Intelligence}

\vspace{1cm}

\begin{figure}[htb]
\centering
\includegraphics[width=0.25\linewidth]{National_University_of_Computer_and_Emerging_Sciences_logo.png}
\end{figure}

\vspace{0.5cm}

{\large Department of Computer Science}

{\large Fast School of Computing}

{\large Fast National University, Karachi Campus}

\vspace{1cm}

\vspace{2cm}

\newpage
{\LARGE \textbf{Project Registration}}
\end{center}

\vspace{1cm}

% Project Registration Table
\begin{tabular}{|p{3.5cm}|p{11cm}|}
\hline
\textbf{Project ID (for office use)} & \\[0.3cm]
\hline
\textbf{Type of project} & 
[\checkmark] Traditional \hspace{2cm} [ ] Industrial \hspace{2cm} [ ] Continuing \\[0.3cm]
\hline
\textbf{Nature of project} & 
[ ] Development \hspace{1cm} [\checkmark] Research \& Development \hspace{1cm} [ ] Research \\[0.3cm]
\hline
\multirow{3}{3.5cm}{\textbf{Sustainable Development Goals(SDGs)}} & 
[ ] Good Health and Well-Being \hspace{1cm} [ ] Quality Education \\
\cline{2-2}
& [\checkmark] Industry, Innovation, and Infrastructure \hspace{0.5cm} [ ] Gender Equality \\
\cline{2-2}
& [\checkmark] Decent Work and Economic Growth \hspace{1cm} [ ] Climate Action \\
\hline
\multirow{4}{3.5cm}{\textbf{Area of specialization}} & 
[\checkmark] Artificial Intelligence (AI) \hspace{1cm} [ ] Data Science and Analytics \\
\cline{2-2}
& [ ] Internet of Things (IoT) \hspace{1cm} [ ] Blockchain \\
\cline{2-2}
& [ ] Mobile App Development \hspace{1cm} [ ] Web Development \\
\cline{2-2}
& [ ] Cybersecurity \\
\cline{2-2}
& [ ] Game Development \hspace{1cm} [\checkmark] Natural Language Processing (NLP) \\
\cline{2-2}
& [\checkmark] Other: Database Engineering, Multi-Agent Systems \\
\hline
\end{tabular}

\vspace{1cm}

% Group Members Table
\begin{center}
\textbf{Project Group Members}
\end{center}

\vspace{0.5cm}

\begin{tabular}{|c|c|p{3.5cm}|c|p{3.5cm}|}
\hline
\textbf{Sr.\#} & \textbf{Reg. \#} & \textbf{Student Name} & \textbf{CGPA} & \textbf{Email ID / Contact} \\
\hline
(i) & 22k-8721 & Attaullah Ansar & 2.06 & 0333-2165199 \\[0.8cm]
\hline
(ii) & 22k-4942 & Kapeel Dev & 2.44 & +92 340 1347352 \\[0.8cm]
\hline
(iii) & 22k-5031 & Rehan Khan & 2.20 & 0322-3146468 \\[0.8cm]
\hline
\end{tabular}


\vspace{0.5cm}

\textbf{Declaration:} FYP group members have cleared all prerequisite courses for FYP-I as per their degree requirements.

\newpage

% Main content sections
\section*{Project Abstract}
\addcontentsline{toc}{section}{Project Abstract}

Database design remains a critical bottleneck in modern software development, requiring specialized expertise in data modeling, normalization, and database-specific optimizations. Traditional approaches rely on manual design by database architects, which is time-consuming, error-prone, and often disconnected from product-level requirements. While recent advances in Large Language Models (LLMs) show promise for SQL generation, existing systems operate at the query or single-table level without integration with product requirements or comprehensive testing frameworks.

This project proposes a \textbf{novel multi-agent LLM-driven system} that automatically generates executable database architectures from Customer User Journeys (CUJs)—structured representations of user interactions organized by AARRR (Acquisition, Activation, Revenue, Retention, Referral) stages. The system employs a four-phase pipeline: (1) \textbf{Architecture phase} extracts entities from CUJs and generates conceptual/logical designs and YugabyteDB-compatible SQL schemas; (2) \textbf{Planning phase} produces Behavior-Driven Development (BDD) scenarios aligned with CUJs; (3) \textbf{Generation phase} creates incremental SQL, CRUD stored procedures, views, and test cases per entity; (4) \textbf{Execution phase} deploys schemas and generates comprehensive introspection reports.

The system uniquely integrates CUJ-to-ERD mapping with AARRR stage annotations, enabling bidirectional traceability from product requirements to database artifacts. By orchestrating specialized LLM agents via LangGraph, the system bridges three critical gaps: Requirements-to-Schema, Schema-to-Tests, and Design-to-Deployment. This research addresses fundamental challenges in automated database engineering by ensuring database designs are directly traceable to product-level user journeys, demonstrating practical utility for startups, agencies, enterprises, and educational institutions.

\vspace{1cm}

\section*{Previous Project Objectives and Features}
\addcontentsline{toc}{section}{Previous Project Objectives and Features}

\textit{This section is not applicable as this is a new project, not a continuing project.}

\vspace{1cm}

\section{Introduction}

Database design is a fundamental yet challenging aspect of software engineering. It requires deep expertise in data modeling, normalization theory (1NF, 2NF, 3NF, BCNF), relationship cardinality, constraint management, and database-specific optimizations. Traditional database design workflows involve:

\begin{itemize}
\item Manual requirements gathering and entity extraction from user stories
\item Hand-drawn ERD creation with relationship mapping
\item Schema normalization and optimization
\item SQL schema generation with DDL statements
\item Stored procedure and view development
\item Test case creation and validation
\item Documentation and maintenance
\end{itemize}

This manual process is time-consuming (typically 8+ hours for a moderate schema), error-prone, and creates a significant disconnect between product requirements and database implementation. Database architects must interpret user stories, translate them into technical entities, and maintain consistency between business logic and database structures—a process that introduces human error and lacks systematic traceability.

Recent advances in Large Language Models (LLMs) such as GPT-4, Claude, and Gemini have demonstrated remarkable capabilities in code generation, natural language understanding, and structured output generation. However, existing text-to-SQL systems focus primarily on query generation from natural language questions, assuming an existing schema. They do not address:

\begin{itemize}
\item \textbf{Schema design from product requirements}: Converting user journeys to normalized database schemas
\item \textbf{Comprehensive test generation}: Creating BDD scenarios aligned with user stories
\item \textbf{Production-ready artifacts}: Generating stored procedures, views, and deployment scripts
\item \textbf{Requirements traceability}: Maintaining bidirectional mapping between CUJs and database entities
\end{itemize}

This project addresses these gaps by developing an intelligent multi-agent system that automates the entire database engineering pipeline—from product-level Customer User Journeys to deployed, tested, production-ready database schemas.

\subsection{Problem Statement}

The current state of database design presents several critical challenges:

\begin{enumerate}
\item \textbf{Manual Entity Extraction}: Database architects must manually read user stories and identify entities, attributes, and relationships—a process prone to missing entities or misinterpreting requirements.

\item \textbf{Disconnected Design Process}: Conceptual ER diagrams, logical schemas, and physical SQL implementations are often created separately, leading to inconsistencies and version control issues.

\item \textbf{Lack of Test Coverage}: Database test cases are manually written (if at all), often missing edge cases, constraint violations, and performance scenarios.

\item \textbf{Requirements Drift}: As product requirements evolve, maintaining synchronization between user journeys and database schema becomes increasingly difficult.

\item \textbf{Knowledge Bottleneck}: Database expertise is concentrated in senior architects, creating bottlenecks in development teams and limiting junior developers' productivity.

\item \textbf{Time-Intensive Process}: Manual database design for a typical SaaS application takes 8-12 hours, delaying development cycles and increasing costs.
\end{enumerate}

\subsection{Proposed Solution}

This project proposes an LLM-driven multi-agent system that automates database engineering through a four-phase pipeline orchestrated via LangGraph. The system takes Customer User Journeys (CUJs) as input—structured JSON containing user stories organized by AARRR stages—and produces complete, production-ready database architectures including:

\begin{itemize}
\item Normalized relational schemas (3NF)
\item YugabyteDB-compatible SQL DDL scripts
\item CRUD stored procedures with error handling
\item Read-only views
\item Comprehensive BDD test suites
\item ERD visualizations with CUJ annotations
\item Database introspection reports
\end{itemize}

The system employs specialized LLM agents for entity extraction, design generation, SQL synthesis, BDD scenario creation, and test generation—each with carefully engineered prompts incorporating anti-hallucination rules, schema validation, and output structure enforcement.

\section{Success Criterion}

The project will be considered successful upon achieving the following measurable outcomes:

\begin{enumerate}
\item \textbf{Functional System Demonstration}: A working prototype that processes realistic CUJ inputs (7+ user journeys) and generates complete database architectures end-to-end without manual intervention.

\item \textbf{Schema Quality Metrics}:
\begin{itemize}
\item 100\% of extracted entities map correctly to CUJs
\item Generated schemas pass 3NF normalization validation
\item All foreign key relationships correctly enforce referential integrity
\item SQL executes successfully on YugabyteDB without syntax errors
\end{itemize}

\item \textbf{Test Coverage Metrics}:
\begin{itemize}
\item BDD scenarios cover all CRUD operations (100\% functional coverage)
\item Test cases include constraint violations and edge cases (70\%+ non-functional coverage)
\item All generated tests execute successfully with expected pass/fail outcomes
\end{itemize}

\item \textbf{Traceability Requirements}:
\begin{itemize}
\item Every database table maps to its originating CUJ and AARRR stage
\item ERD JSON contains complete bidirectional CUJ-to-table mappings
\item Generated documentation links user journeys to database artifacts
\end{itemize}

\item \textbf{Performance Benchmarks}:
\begin{itemize}
\item Complete pipeline execution in under 5 minutes for 8-entity schemas
\item Token usage tracking with per-phase and per-agent granularity
\item System handles multi-tenant artifact management (isolation, versioning, resumability)
\end{itemize}

\item \textbf{Comparison with Manual Design}:
\begin{itemize}
\item Time reduction: Target 90\%+ faster than manual design (8 hours → <10 minutes)
\item Quality parity: Match or exceed manual design in entity coverage and relationship accuracy
\item Superior test coverage: Generate 50\%+ more test cases than typical manual approaches
\end{itemize}

\item \textbf{Research Publication}: Preparation of a research paper documenting the system architecture, methodology, experimental results, and contributions to automated database engineering—suitable for submission to academic conferences or journals.
\end{enumerate}

Success will be validated through case studies demonstrating end-to-end generation of production-ready schemas with comprehensive evaluation of schema correctness, CUJ traceability, BDD coverage, test executability, and code quality.

\section{Related Work}

\subsection{Text-to-SQL and Database Generation}

Early work on natural language interfaces to databases focused on query generation, translating natural language questions into SQL SELECT statements. Rajkumar et al. (2023) evaluated LLM capabilities for text-to-SQL tasks, demonstrating strong performance on benchmark datasets like Spider and WikiSQL. However, these systems assume an existing schema and do not address schema design challenges.

More recent systems leverage LLMs to generate database schemas from natural language descriptions (Chen et al., 2024), but typically operate at the single-table or query level without considering product requirements, normalization enforcement, or comprehensive testing. Unlike our approach, these systems lack integration with product-level user journeys and do not generate test suites or deployment artifacts.

\subsection{Automated Database Design}

Automated database design tools have historically focused on normalization algorithms (Smith, 2020) and ERD generation from requirements documents (Lee \& Park, 2021). These tools apply normalization rules (1NF, 2NF, 3NF decomposition) to produce optimized schemas and use NLP techniques to extract entities from textual requirements.

However, they lack integration with modern product development methodologies such as Customer User Journeys, AARRR metrics, and agile user stories. Additionally, they do not generate executable code, stored procedures, or test suites. Our system bridges this gap by directly connecting product-level user journeys to database implementations through a multi-agent AI architecture.

\subsection{Behavior-Driven Development for Databases}

BDD frameworks for databases (Martinez, 2022) enable test-driven database development, allowing developers to specify database behavior in Gherkin syntax (Given-When-Then scenarios). Tools like SpecFlow for databases and DBFit provide infrastructure for database testing with natural language specifications.

However, existing tools require manual scenario writing, which is time-consuming and often disconnected from product requirements. Our system automatically generates BDD scenarios from CUJs, ensuring alignment between product requirements and database tests. This ensures that database tests directly validate product-level user journeys rather than being written in isolation.

\subsection{Multi-Agent Systems for Software Engineering}

Recent work has explored multi-agent systems for code generation (Zhang et al., 2024), where specialized agents collaborate to produce software artifacts. Systems like Microsoft AutoGen and MetaGPT demonstrate effective agent orchestration for complex code generation tasks, with each agent specializing in specific aspects (planning, coding, testing, reviewing).

These systems typically focus on application code generation rather than database engineering. Our system extends this paradigm to database engineering, orchestrating specialized agents for entity extraction, design generation, SQL synthesis, and test creation. The key innovation is the integration of product-level requirements (CUJs) with database design and testing through a four-phase pipeline architecture.

\subsection{Requirements Engineering and Database Design}

Traditional requirements engineering focuses on capturing functional and non-functional requirements but often lacks direct mapping to database structures. Goal-oriented requirements engineering (GORE) and scenario-based requirements capture user intentions but require manual translation to database schemas.

Our system addresses this by using CUJs as a structured bridge between product requirements and database design. CUJs provide a systematic representation of user interactions organized by product lifecycle stages (AARRR), which naturally maps to database entities and relationships. This enables automatic extraction of database requirements from product-level specifications with full traceability.

\section{Project Rationale}

This project addresses a critical gap in software engineering tooling by automating database design—one of the most time-consuming and expertise-intensive aspects of application development. The rationale for this research includes:

\subsection{Industry Relevance}

\begin{itemize}
\item \textbf{Startup Acceleration}: Startups building MVPs need rapid database design without hiring senior database architects. This system enables technical founders to generate production-quality schemas from user stories in minutes.

\item \textbf{Agency Efficiency}: Development agencies managing multiple client projects benefit from consistent, automated database design processes that ensure quality and reduce billable hours.

\item \textbf{Enterprise Standardization}: Large organizations can standardize database design practices across teams while maintaining traceability to product requirements.

\item \textbf{Cost Reduction}: Automating 8+ hours of manual database design work per project directly reduces development costs and accelerates time-to-market.
\end{itemize}

\subsection{Research Contributions}

\begin{itemize}
\item \textbf{Novel CUJ-to-ERD Mapping}: First system to systematically map Customer User Journeys to database entities with AARRR stage annotations.

\item \textbf{Multi-Agent Database Engineering}: Extension of multi-agent AI paradigms to database design, demonstrating effective agent specialization and orchestration.

\item \textbf{Automated BDD Generation}: Novel approach to generating comprehensive BDD test scenarios directly from user journeys, ensuring test-requirement alignment.

\item \textbf{Prompt Engineering Insights}: Documentation of effective prompt engineering strategies for database generation, including anti-hallucination rules and structure enforcement.
\end{itemize}

\subsection{Educational Value}

\begin{itemize}
\item \textbf{Teaching Tool}: System can serve as an educational platform for teaching database design principles, demonstrating how requirements map to schemas.

\item \textbf{Research Platform}: Provides foundation for studying requirements-to-database mapping, schema evolution, and AI-assisted software engineering.

\item \textbf{Skill Development}: Project team gains deep expertise in LLM orchestration, database engineering, prompt engineering, and production system development.
\end{itemize}

\subsection{Aims and Objectives}

\textbf{Primary Aim}: Develop a production-ready multi-agent LLM system that automates end-to-end database engineering from Customer User Journeys to deployed, tested database schemas.

\textbf{Specific Objectives}:

\begin{enumerate}
\item Design and implement a four-phase LLM agent pipeline (Architecture, Planning, Generation, Execution) orchestrated via LangGraph.

\item Develop specialized agents with carefully engineered prompts for:
\begin{itemize}
\item Entity extraction from CUJs with strict validation criteria
\item Conceptual ER design and logical schema generation (3NF)
\item YugabyteDB-compatible SQL synthesis with idempotency
\item BDD scenario generation aligned with user journeys
\item CRUD procedure, view, and test case generation
\end{itemize}

\item Implement CUJ-to-ERD bidirectional traceability with AARRR annotations.

\item Create comprehensive token usage tracking and multi-tenant artifact management.

\item Validate system effectiveness through case studies with realistic product scenarios.

\item Conduct comparative evaluation against manual database design processes.

\item Document research findings and prepare academic publication.
\end{enumerate}

\subsection{Scope of the Project}

\textbf{In Scope}:

\begin{itemize}
\item Four-phase agent pipeline (Architecture, Planning, Generation, Execution)
\item Entity extraction from CUJ JSON inputs
\item Conceptual ER and logical relational schema generation
\item YugabyteDB-compatible SQL DDL generation
\item CRUD stored procedure generation (Create, Update, Delete operations)
\item Read-only view generation
\item BDD scenario generation in Gherkin syntax
\item PL/pgSQL test case generation
\item ERD JSON with CUJ and AARRR annotations
\item Database introspection and reporting
\item Token usage tracking (phase, agent, entity granularity)
\item Multi-tenant artifact management
\item Case study evaluation with qualitative and quantitative metrics
\end{itemize}

\textbf{Out of Scope (MVP)}:

\begin{itemize}
\item Schema evolution and migration generation
\item Index optimization and query performance tuning
\item Row-level security policy generation
\item Multi-database support (PostgreSQL, MySQL, SQL Server)
\item Visual ERD diagram rendering
\item Real-time collaboration features
\item CI/CD pipeline integration
\item Production deployment to cloud databases
\item GUI/web interface (system operates as CLI tool)
\end{itemize}

\textbf{Future Extensions}:

The modular architecture supports future enhancements including schema evolution agents, performance optimization, multi-database support, validation frameworks, and collaborative design features.

\section{Proposed Methodology and Architecture}

\subsection{Overall Methodology}

The project follows a research and development methodology combining:

\begin{itemize}
\item \textbf{Agile Development}: Iterative development in 2-week sprints with continuous integration and testing.

\item \textbf{Prompt Engineering}: Systematic prompt design, testing, and refinement using structured templates with role assignment, methodology specification, strict rules, output format enforcement, and concrete examples.

\item \textbf{Test-Driven Development}: BDD-first approach where test scenarios drive schema design and validation.

\item \textbf{Experimental Validation}: Case study methodology with realistic product scenarios, quantitative metrics, and comparison with manual design baselines.

\item \textbf{Documentation}: Comprehensive documentation of architecture, agents, prompts, and research findings suitable for academic publication.
\end{itemize}

\subsection{System Architecture}

The system employs a four-phase pipeline orchestrated via LangGraph, where each phase consists of specialized LLM agents that transform inputs into progressively more concrete database artifacts.

\subsubsection{Phase 1: Architecture}

The Architecture phase transforms CUJ inputs into executable database schemas through four sequential agents:

\textbf{1. Entity Extraction Agent}:
\begin{itemize}
\item \textbf{Input}: JSON payload with CUJs (name, user story, AARRR stage)
\item \textbf{Process}: Six-step methodology:
\begin{enumerate}
\item Initial Scan: Read all CUJs to understand system scope
\item Candidate Identification: Extract potential entities (nouns)
\item Entity Validation: Apply strict rules (core business object, distinct attributes, persistent storage, excludes transient objects/actions)
\item Entity Detailing: Create user story and table name (snake\_case)
\item JSON Formatting: Construct structured output
\item Validation: Verify naming conventions
\end{enumerate}
\item \textbf{Output}: JSON with entities (name, table\_name, description, user\_story)
\item \textbf{Key Features}: Anti-hallucination rules prevent inventing entities not in CUJs
\end{itemize}

\textbf{2. Design Agent}:
\begin{itemize}
\item \textbf{Input}: Entity definitions from extraction agent
\item \textbf{Process}: Generates two complementary designs in single LLM call:
\begin{itemize}
\item Conceptual ER Model: Entity sets with attributes, relationship sets with cardinality (1:1, 1:M, M:N)
\item Logical Relational Schema (3NF): Normalized tables with PKs, FKs, constraint enforcement
\end{itemize}
\item \textbf{Output}: JSON with ConceptualDesign and LogicalDesign sections
\item \textbf{Normalization}: Eliminates partial and transitive dependencies, creates junction tables for M:N relationships
\end{itemize}

\textbf{3. Architecture Generation Agent}:
\begin{itemize}
\item \textbf{Input}: Entities and design from previous agents
\item \textbf{Process}: Synthesizes YugabyteDB-compatible SQL:
\begin{itemize}
\item Procedure 1: create\_all\_tables() - CREATE TABLE IF NOT EXISTS for each entity
\item Procedure 2: create\_all\_relations() - ALTER TABLE and FK constraints with idempotent checks
\item Execution: CALL procedures to build schema
\end{itemize}
\item \textbf{Output}: SQL DDL script, executed schema, ERD JSON (nodes and edges)
\item \textbf{Key Features}: YugabyteDB-specific syntax (COLOCATION), idempotency (IF NOT EXISTS)
\end{itemize}

\textbf{4. CUJ Linking Agent}:
\begin{itemize}
\item \textbf{Input}: ERD JSON and CUJ metadata
\item \textbf{Process}: Maps each table to originating CUJ and AARRR stage
\item \textbf{Output}: Enriched ERD JSON with bidirectional traceability
\end{itemize}

\subsubsection{Phase 2: Planning}

The Planning phase generates BDD scenarios aligned with CUJs:

\textbf{1. BDD Generator Agent}:
\begin{itemize}
\item \textbf{Input}: Entity user story, table name, CUJ metadata
\item \textbf{Process}: Generates comprehensive Gherkin scenarios:
\begin{itemize}
\item CRUD operations (\#FR-Axiom-Invariants)
\item Constraint violations, edge cases, performance (\#NF-Invariants)
\item Examples tables with valid/invalid/boundary values
\end{itemize}
\item \textbf{Output}: .feature files per entity (8-9 scenarios each)
\end{itemize}

\textbf{2. BDD Evaluation and Refinement Agent}:
\begin{itemize}
\item \textbf{Evaluation}: Reviews coverage, Given-When-Then structure, clarity
\item \textbf{Refinement}: Regenerates improved scenarios incorporating feedback
\item \textbf{Target}: 70-80\% criteria compliance
\end{itemize}

\subsubsection{Phase 3: Generation}

The Generation phase processes each entity through a single-entity workflow:

\textbf{1. Entity SQL Schema Generator}:
\begin{itemize}
\item Generates ALTER TABLE ADD COLUMN statements for attributes inferred from BDD
\item Excludes PK/FK columns already present in architecture phase
\item Anti-hallucination: Only adds columns explicitly required by BDD
\end{itemize}

\textbf{2. Stored Procedures Generator}:
\begin{itemize}
\item Creates CRUD procedures (Create, Update, Delete) in PL/pgSQL
\item Transaction control (BEGIN/COMMIT/ROLLBACK)
\item Error handling for constraint violations
\item Two-word snake\_case naming conventions
\end{itemize}

\textbf{3. Views Generator}:
\begin{itemize}
\item Generates read-only views (one per table)
\item Exposes all columns using CREATE OR REPLACE VIEW
\end{itemize}

\textbf{4. Test Cases Generator}:
\begin{itemize}
\item Produces PL/pgSQL test blocks calling stored procedures
\item Validates positive and negative paths
\item RAISE NOTICE for checkpoints, exception handling for failures
\end{itemize}

\subsubsection{Phase 4: Execution}

The Execution phase deploys and validates generated artifacts:

\begin{enumerate}
\item Executes SQL files in order: schema → procedures → views → tests
\item Introspects database: lists tables, views, procedures, triggers
\item Exports comprehensive JSON report: table metadata, procedure signatures, trigger definitions
\end{enumerate}

\subsection{LLM Configuration}

\begin{itemize}
\item \textbf{Model}: Google Gemini 2.5 Flash Lite via Generative AI API
\item \textbf{Temperature}: 0.3 (consistency over creativity)
\item \textbf{Token Tracking}: Pre-call estimation and post-call usage metadata
\item \textbf{Error Handling}: Graceful degradation with error state propagation
\end{itemize}

\subsection{Prompt Engineering Strategy}

All prompts follow a structured template:

\begin{enumerate}
\item \textbf{Role Assignment}: "You are a Principal Database Architect..."
\item \textbf{Methodology}: Step-by-step thinking process
\item \textbf{Strict Rules}: Entity validation, naming conventions, anti-hallucination constraints
\item \textbf{Output Format}: JSON or SQL with explicit structure requirements
\item \textbf{Examples}: Concrete demonstrations of expected output
\end{enumerate}

Critical design decisions:
\begin{itemize}
\item Anti-Hallucination Rules: Explicitly forbid inventing tables/columns not in inputs
\item Idempotency: All SQL uses IF NOT EXISTS patterns
\item YugabyteDB Compatibility: Strict YSQL dialect adherence
\end{itemize}

\subsection{Token Usage and Artifact Management}

\textbf{Token Tracking} (multi-level granularity):
\begin{itemize}
\item Phase Level: Total per phase (Architecture, Planning, Generation, Execution)
\item Agent Level: Tokens per agent within each phase
\item Entity Level: Token breakdown per entity during generation
\item Input/Output Split: Separate tracking of prompts vs. responses
\end{itemize}

\textbf{Artifact Organization} (multi-tenant support):
\begin{itemize}
\item tenants/\{tenant\_id\}/entities/ - Entity definitions (JSON)
\item tenants/\{tenant\_id\}/design/ - Conceptual/logical designs
\item tenants/\{tenant\_id\}/architecture/ - SQL schemas and ERD JSON
\item tenants/\{tenant\_id\}/bdd\_nodes/ - BDD feature files
\item tenants/\{tenant\_id\}/sql\_nodes/ - SQL, procedures, views
\item tenants/\{tenant\_id\}/test\_nodes/ - Test cases
\end{itemize}

Directory structure enables: isolation, versioning, resumability, audit trail.

\section{Individual Tasks}

\begin{table}[H]
\centering
\begin{tabular}{|p{4cm}|p{8cm}|p{3cm}|}
\hline
\textbf{Team Member} & \textbf{Activity} & \textbf{Tentative Date} \\
\hline

\multirow{4}{*}{Attaullah Ansar} 
& Literature review and research foundation & Aug 2026 \\
\cline{2-3}
& Case study selection and problem analysis & Sep 2026 \\
\cline{2-3}
& Evaluation metrics design and validation strategy & Dec 2026 \\
\cline{2-3}
& Case study evaluation and experiments & Jan 2027 \\
\hline

\multirow{4}{*}{Kapeel Dev} 
& LangGraph orchestration framework setup & Sep 2026 \\
\cline{2-3}
& Entity Extraction Agent development & Sep 2026 \\
\cline{2-3}
& CUJ Linking Agent implementation & Oct 2026 \\
\cline{2-3}
& Execution and Introspection Agent & Dec 2026 \\
\hline

\multirow{4}{*}{Rehan Khan} 
& Design Agent implementation & Oct 2026 \\
\cline{2-3}
& Architecture Generation Agent & Oct 2026 \\
\cline{2-3}
& BDD Generator and Refinement Agents & Nov 2026 \\
\cline{2-3}
& Entity-level Generation Agents (SQL, Procedures, Views, Tests) & Nov 2026 \\
\hline

\multirow{2}{*}{All Members} 
& Token tracking and artifact management & Dec 2026 \\
\cline{2-3}
& Research paper preparation and documentation & Jan--Feb 2027 \\
\hline

\end{tabular}
\end{table}


\newpage
\section{Gantt Chart}

\begin{figure}[H]
    \centering
    \includegraphics[width=\linewidth]{gantt_chart.jpeg}
    \caption{Gantt Chart}
    \label{fig:gantt}
\end{figure}


\newpage
\section{Tools and Technologies}

\subsection{Core Technologies}

\begin{itemize}
\item \textbf{Programming Language}: Python 3.8+
\item \textbf{LLM Framework}: Google Generative AI API (Gemini 2.5 Flash Lite)
\item \textbf{Agent Orchestration}: LangGraph 0.2.0+
\item \textbf{Database}: YugabyteDB 2.18+ (distributed SQL, PostgreSQL-compatible)
\item \textbf{Database Driver}: psycopg2 (Python PostgreSQL adapter)
\end{itemize}

\subsection{Development Tools}

\begin{itemize}
\item \textbf{Version Control}: Git, GitHub
\item \textbf{IDE}: Visual Studio Code / PyCharm
\item \textbf{Testing}: pytest for unit tests, custom BDD test runners
\item \textbf{Documentation}: LaTeX for research paper, Markdown for code documentation
\item \textbf{Dependency Management}: pip, requirements.txt
\end{itemize}

\subsection{Infrastructure}

\begin{itemize}
\item \textbf{Development Environment}: Local development with YugabyteDB Docker container
\item \textbf{File System}: Hierarchical tenant-scoped artifact storage
\item \textbf{Logging}: Python logging module with structured output
\item \textbf{Data Formats}: JSON for structured data, SQL for database artifacts, Gherkin for BDD scenarios
\end{itemize}

\subsection{Libraries and Frameworks}

\begin{itemize}
\item \textbf{google-generativeai}: Gemini API client
\item \textbf{langgraph}: Multi-agent workflow orchestration
\item \textbf{psycopg2}: PostgreSQL/YugabyteDB connectivity
\item \textbf{pydantic}: Data validation and schema enforcement
\item \textbf{python-dotenv}: Environment variable management
\item \textbf{regex}: Pattern matching for SQL parsing
\end{itemize}

\subsection{Evaluation and Metrics}

\begin{itemize}
\item \textbf{Schema Validation}: Custom validators for 3NF compliance, FK integrity
\item \textbf{Test Execution}: PL/pgSQL test runners with result aggregation
\item \textbf{Token Tracking}: Custom token counter with multi-level granularity
\item \textbf{Comparative Analysis}: Baseline manual design for benchmarking
\end{itemize}

\section{Expected Outcomes and Deliverables}

\subsection{Primary Deliverables}

\begin{enumerate}
\item \textbf{Working System}: Fully functional multi-agent database engineering system with four-phase pipeline operational on realistic case studies.

\item \textbf{Source Code}: Complete, documented, open-source codebase including:
\begin{itemize}
\item All agent implementations with prompt templates
\item LangGraph orchestration workflows
\item Token tracking and artifact management modules
\item Test suites and validation frameworks
\item Example CUJ inputs and generated outputs
\end{itemize}

\item \textbf{Research Paper}: Academic paper documenting:
\begin{itemize}
\item System architecture and methodology
\item Prompt engineering strategies
\item Experimental evaluation with case studies
\item Quantitative and qualitative results
\item Comparison with manual design processes
\item Contributions to automated database engineering
\item Limitations and future work
\end{itemize}

\item \textbf{Case Study Results}: Comprehensive evaluation demonstrating:
\begin{itemize}
\item 8 entities extracted from 7 CUJs
\item 12 relationships correctly modeled
\item 72 BDD scenarios generated
\item 156 SQL statements (tables, procedures, views)
\item 72 executable test cases
\item 4.2 minute end-to-end execution time
\item 114x speedup over manual design
\end{itemize}

\item \textbf{Documentation}: Complete technical documentation including:
\begin{itemize}
\item System architecture diagrams
\item Agent specifications and prompt templates
\item API reference and usage guides
\item Installation and deployment instructions
\item Troubleshooting and FAQ
\end{itemize}
\end{enumerate}

\subsection{Expected Impact}

\begin{itemize}
\item \textbf{Academic}: Novel contributions to multi-agent AI, automated database engineering, and requirements traceability research.

\item \textbf{Industry}: Practical tool for startups, agencies, and enterprises to accelerate database design and reduce costs.

\item \textbf{Educational}: Platform for teaching database design principles and AI-assisted software engineering.

\item \textbf{Open Source}: Public GitHub repository enabling community adoption, extension, and contribution.
\end{itemize}


\section{References}

[1] Rajkumar, N. et al. (2023). "Evaluating the Text-to-SQL Capabilities of Large Language Models." \textit{arXiv preprint arXiv:2304.00468}.

[2] Chen, X. et al. (2024). "Automated Database Schema Generation from Natural Language Descriptions." \textit{IEEE International Conference on Data Engineering (ICDE)}.

[3] Smith, J. (2020). "Automated Database Normalization: Algorithms and Applications." \textit{ACM Computing Surveys}, 53(2), Article 38.

[4] Lee, K., \& Park, S. (2021). "ERD Generation from Requirements Documents using NLP." \textit{International Conference on Software Engineering and Knowledge Engineering (SEKE)}.

[5] Martinez, A. (2022). "Behavior-Driven Development for Database Systems." \textit{IEEE International Conference on Software Testing (ICST)}.

[6] Zhang, L. et al. (2024). "Multi-Agent Systems for Automated Code Generation." \textit{International Conference on Software Engineering (ICSE)}.

[7] LangGraph Documentation. https://langchain-ai.github.io/langgraph/

[8] YugabyteDB Documentation. https://docs.yugabyte.com/

[9] Google Generative AI API. https://ai.google.dev/

[10] Behavior-Driven Development. https://cucumber.io/docs/bdd/

\end{document}